\documentclass[11pt]{article}
\usepackage[margin=1.15in]{geometry}
\usepackage{amsmath,amssymb,amsthm,mathtools}
\usepackage{booktabs}
\usepackage[hidelinks]{hyperref}

\theoremstyle{definition}
\newtheorem{definition}{Definition}[section]
\theoremstyle{plain}
\newtheorem{theorem}[definition]{Theorem}
\newtheorem{proposition}[definition]{Proposition}
\newtheorem{lemma}[definition]{Lemma}
\newtheorem{corollary}[definition]{Corollary}
\newtheorem{conjecture}[definition]{Conjecture}
\theoremstyle{remark}
\newtheorem{measurement}[definition]{Measurement}
\newtheorem{remark}[definition]{Remark}

\newcommand{\Adm}{\mathrm{Adm}}
\newcommand{\Fix}{\operatorname{Fix}}
\newcommand{\lfp}{\operatorname{lfp}}
\newcommand{\gfp}{\operatorname{gfp}}
\newcommand{\El}{\mathcal{E}}
\newcommand{\Law}{\mathcal{R}}
\newcommand{\Haz}{\mathcal{H}}
\newcommand{\War}{\mathcal{W}}

\title{An Algebra of Mechanism Composition\\
\large Structure and obstruction in a vocabulary of on-chain financial mechanisms}
\author{}
\date{}

\begin{document}
\maketitle

\begin{abstract}
We study a finite vocabulary $\El$ of recurring on-chain financial mechanisms
together with two families of constraints: \emph{requirements}, stating what a
mechanism needs, and \emph{prohibitions}, stating which combinations are
forbidden. We show that admissibility is the set of common fixed points of an
inflationary operator $\Gamma$ and a deflationary operator $\Delta$; that this
``diagonal'' structure obstructs three standard frameworks (interlaced
bilattices, approximation fixpoint theory, and equalizers in $\mathbf{Pos}$);
that reasoning over the full powerset collapses for reasons intrinsic to the
constraint set rather than to any operator; and that a sound and empirically
exact procedure exists on \emph{completion lattices} $[S,\top]$. We record what
is proved, what is measured, and what remains open, and we are explicit about
which is which.
\end{abstract}

\section*{Status conventions}

This document is maintained incrementally. Every numbered item carries its
epistemic status in its environment, and the distinction is load-bearing:

\begin{itemize}
  \item \textbf{Theorem} / \textbf{Proposition} / \textbf{Corollary} --- proved,
    with the proof given or sketched. A sketch means the argument is complete and
    the write-up is not.
  \item \textbf{Measurement} --- established computationally against the corpus
    or by exhaustive enumeration, with the instance and bound stated. Not proved
    in general, and \emph{not} to be cited as if it were.
  \item \textbf{Conjecture} --- believed, with the evidence named. Several items
    here began as measurements and are awaiting proof; several began as claims we
    made and were subsequently refuted.
\end{itemize}

Where a result of ours was later overturned, the overturning is recorded rather
than the earlier claim silently removed. The running status of every item is
tracked alongside this document in \texttt{algebra/THEOREM-LEDGER.md}.

\section{Preliminaries}

\begin{definition}[Vocabulary and protocols]
Let $\El$ be a finite set of \emph{elements} (mechanisms), $|\El| = 58$. A
\emph{protocol} is a subset $X \subseteq \El$. We write
$L = (2^{\El}, \subseteq)$ for the resulting complete lattice, with
$\bot = \emptyset$ and $\top = \El$.
\end{definition}

\begin{definition}[Requirements]
A \emph{requirement} is a pair $(s, T)$ with $s \in \El$ a \emph{subject} and
$T = \{T_1,\dots,T_k\}$ a finite family of \emph{terms}, each $T_j \subseteq \El$
a nonempty set of \emph{alternatives}. It is \emph{satisfied} by $X$ when
\[
  s \in X \;\Longrightarrow\; \forall j \le k,\; T_j \cap X \neq \emptyset .
\]
Write $\Law$ for the set of requirements. In clause form a requirement is
$\neg s \vee \bigvee_{e \in T_j} e$, which carries exactly one negative literal
and is therefore \emph{dual-Horn}.
\end{definition}

\begin{definition}[Prohibitions]
A \emph{prohibition} is a set $H \subseteq \El$, forbidden in the sense that
$H \not\subseteq X$ is required. Write $\Haz$ for the family of prohibitions; we
take $\Haz$ to be an antichain (a \emph{clutter}) of minimal forbidden sets. In
clause form a prohibition is $\bigvee_{e \in H} \neg e$: purely negative, hence
\emph{Horn}.
\end{definition}

\begin{definition}[Warrants]
A \emph{warrant} is a pair $(e, C)$ with $e \in \El$ and $C \subseteq \El$ a set
of \emph{consumers}. It is satisfied by $X$ when
$e \in X \Rightarrow C \cap X \neq \emptyset$. Write $\War$ for the set of
warrants. Where a requirement says what $e$ \emph{needs}, a warrant says what $e$
is \emph{for}: an element present with no consumer is unwarranted.
\end{definition}

\begin{remark}
$\War$ is not independent data. It is obtained as the residual (order-theoretic
adjoint) of the requirement relation; the vocabulary as originally recorded
contains $\Law$ and $\Haz$ only, and $|\War| = 27$.
\end{remark}

\section{The two operators}

\begin{definition}[Closure and warrant operators]
Define $\Gamma, \Delta : L \to L$ by
\[
  \Gamma(X) \;=\; X \cup \{\,\text{elements forced by requirements fired in } X\,\},
  \qquad
  \Delta(X) \;=\; X \setminus \{\,e \in X : \text{$e$ is unwarranted in } X\,\},
\]
each iterated where necessary to a fixed point of its own defining step.
\end{definition}

\begin{proposition}\label{prop:polarity}
$\Gamma$ is \emph{inflationary} $(X \subseteq \Gamma(X))$, monotone and
idempotent, hence a closure operator. $\Delta$ is \emph{deflationary}
$(\Delta(X) \subseteq X)$ and monotone.
\end{proposition}

\begin{measurement}\label{meas:notidem}
$\Delta$ is \emph{not} idempotent.
\end{measurement}

\begin{proposition}[Free repair]\label{prop:repair}
$\Delta^{\omega} \coloneqq \Delta^{|\El|}$ is deflationary, monotone and
idempotent --- a kernel operator --- and $\Fix(\Delta^{\omega}) = \Fix(\Delta)$.
Consequently Measurement~\ref{meas:notidem} is immaterial to everything below,
and we may take $\Delta$ idempotent without loss.
\end{proposition}

\begin{definition}[Admissibility]
$X$ is \emph{admissible} iff $\Gamma(X) = X$, $\Delta(X) = X$, and no
$H \in \Haz$ satisfies $H \subseteq X$. Write
$\Adm = \Fix(\Gamma) \cap \Fix(\Delta) \cap \Haz^{c}$.
\end{definition}

\begin{proposition}\label{prop:moore}
$\Fix(\Gamma)$ is a Moore family (closed under intersection) and $\Fix(\Delta)$
is an interior system (closed under union). Neither conclusion requires
idempotence of the generating operator.
\end{proposition}

\begin{remark}[The diagonal]
Proposition~\ref{prop:moore} is the structural content of the whole paper:
$\Adm$ is a \emph{Moore family intersected with a co-Moore family}. The two
closure halves cancel exactly, and $\Adm$ is a diagonal --- the agreement set of
two maps pointing in opposite directions relative to the identity. Every
obstruction in Section~\ref{sec:obstructions} is an instance of a framework
being unable to see a diagonal.
\end{remark}

\section{Polarity}

\begin{theorem}[Polarity]\label{thm:polarity}
The models of $\Law$ are closed under union; the models of $\Haz$ are closed
under intersection; the models of $\Law \cup \Haz$ are in general closed under
neither.
\end{theorem}

\begin{proof}[Proof sketch]
Requirement clauses are dual-Horn and prohibition clauses are Horn. By
Schaefer's classification, dual-Horn constraint languages are exactly those
preserved by $\vee$ (union) and Horn languages exactly those preserved by
$\wedge$ (intersection). A conjunction of the two need be preserved by neither,
and explicit witnesses exist in the present instance.
\end{proof}

\begin{measurement}\label{meas:arity}
Of the $34$ element-expressible clauses ($31$ requirement, $3$ prohibition),
exactly $17$ are \emph{bijunctive} (width $\le 2$); the widest has $5$ literals.
Hence the system is not median-closed, though a $50\%$ fragment is.
\end{measurement}

\section{Obstructions}\label{sec:obstructions}

\begin{theorem}[Bilattices]\label{thm:bilattice}
No interlaced bilattice on $\Fix(\Gamma) \times \Fix(\Delta)$ has $\Adm$ as a
definable object.
\end{theorem}

\begin{proof}[Proof sketch]
By Avron's representation theorem every interlaced bilattice is isomorphic to a
componentwise product $L \odot R$, uniquely. Taking $L = \Fix(\Gamma)$ and
$R = \Fix(\Delta)$ yields pairs with no constraint tying the coordinates, whereas
$\Adm$ is precisely the diagonal $\{(x,x)\}$. A componentwise structure cannot
distinguish it. Moreover a bilattice would make $\Adm$ a lattice under both
orders, contradicting Theorem~\ref{thm:polarity}.
\end{proof}

\begin{theorem}[Approximation fixpoint theory]\label{thm:aft}
There is no approximator $A$ on $L^2$ whose lower and upper components are
$\Gamma$ and $\Delta$ respectively, other than the trivial one with
$\Gamma = \Delta = \mathrm{id}$.
\end{theorem}

\begin{proof}
Every approximator is required to preserve consistency: $A^1(x,y) \le A^2(x,y)$.
Evaluating on the diagonal, $A(x,x) = (\Gamma(x), \Delta(x))$ gives
$\Gamma(x) \le \Delta(x)$ for all $x$. But $\Gamma$ is inflationary, so
$x \le \Gamma(x)$, and $\Delta$ is deflationary, so $\Delta(x) \le x$. Chaining,
\[
  \Gamma(x) \;\le\; \Delta(x) \;\le\; x \;\le\; \Gamma(x),
\]
whence $\Gamma(x) = \Delta(x) = x$ for every $x$. The argument uses neither
exactness nor $\le_p$-monotonicity, and applies to any $A$, product-form or not.
\end{proof}

\begin{corollary}
The assignment is backwards: an approximator's lower component must
\emph{under}-approximate, while a closure operator \emph{over}-approximates.
Exchanging the roles is admissible but vacuous, since $\lfp(\Delta) = \bot$
identically, so the induced stable operator is constant.
\end{corollary}

\begin{remark}
A common misdiagnosis should be recorded. It is \emph{not} the case that
monotonicity of both $\Gamma$ and $\Delta$ obstructs $\le_p$-monotonicity of $A$:
antitonicity in the second argument is satisfied vacuously by constancy, and for
monotone $O$ the ultimate approximator is exactly the product
$(O(x), O(y))$. The obstruction is consistency, as above.
\end{remark}

\section{Collapse over the powerset}

\begin{proposition}\label{prop:collapse}
Let $O$ be any operator on $L$ with $\Fix(O) = \Adm$. If
\emph{(i)} $\emptyset \in \Adm$ and \emph{(ii)} $\bigcup \Adm = \El$, then the
ultimate approximator satisfies $U_O(\bot,\top) = (\bot,\top)$; the
Kripke--Kleene fixpoint is $(\bot,\top)$ and carries no information.
\end{proposition}

\begin{proof}
$U_O(x,y) = \bigl(\bigwedge O([x,y]),\ \bigvee O([x,y])\bigr)$. The interval
$[\bot,\top]$ contains $\emptyset$, and by (i) $O(\emptyset) = \emptyset$, so the
lower bound is $\bot$. By (ii) the upper bound is $\top$. Hence
$(\bot,\top)$ is a fixed point of the iteration at the first step.
\end{proof}

\begin{remark}
Hypotheses (i) and (ii) hold for our vocabulary and are facts about the
\emph{constraint set}, not about $O$: (i) because every requirement is an
implication with a nonempty antecedent, so the empty protocol violates nothing;
(ii) because every element occurs in some admissible set. Proposition
\ref{prop:collapse} therefore rules out bottom-up reasoning over $2^{\El}$ for
\emph{every} operator with the right fixpoints. The obstruction is the lattice.
\end{remark}

\section{Completion lattices}

\begin{definition}[Completion lattice]
For a seed $S \subseteq \El$, the \emph{completion lattice} is the interval
$[S,\top] = \{X : S \subseteq X \subseteq \El\}$, and the associated question is
\emph{which extensions of $S$ are admissible}.
\end{definition}

\begin{proposition}
$\Gamma$ restricts to $[S,\top]$, with least fixed point $\Gamma(S)$ reached in
one step. $\Delta$ does not restrict; the corrected operator
$\Delta_S(X) \coloneqq \Delta(X) \cup S$ does.
\end{proposition}

\begin{definition}[Downward iteration]
Set $x_0 = \top$ and $x_{n+1} = \Gamma\bigl(\Delta(x_n) \cup S\bigr)$, and let
$x_\infty$ be the limit.
\end{definition}

\begin{proposition}[The iteration is vacuous here]\label{prop:vacuous}
If every dependent element has a consumer present in $\top$, then
$\Delta(\top) = \top$, and since $\Gamma$ is inflationary the iteration is
stationary at $x_0$: $x_\infty = \top$ for every seed.
\end{proposition}

\begin{proof}
Warrant is a \emph{co-presence} condition: $e$ is unwarranted in $X$ only if
$C(e) \cap X = \emptyset$. At $X = \top$ every consumer set is met, so $\Delta$
removes nothing. Then $x_1 = \Gamma(\top \cup S) = \top$.
\end{proof}

\begin{measurement}\label{meas:vacuous}
On synthetic operators of the present shape the iteration was sound $2000/2000$
and exact $2000/2000$ with zero slack. On the actual $\Gamma,\Delta$ it is sound
and \emph{vacuous}: $x_\infty = \top$ on every seed at $10$, $20$ and $58$
elements, exact in $0$ of $18$ seed-instances, slack $1$--$39$. The synthetic
operators admitted a $\Delta$ that removes at $\top$; ours structurally cannot.
\end{measurement}

\begin{remark}
A ban-aware variant does descend and recovers the atomic-scope prohibition on
seed $\{Fl,Xm\}$, but it is \emph{unsound}: on seed $\{Uc\}$ it prunes to
$\{Uc,Aw,At\}$, excluding every witness of the third and fourth terms of the
governing requirement, and so declares $Uc$ uncompletable although two live
protocols instantiate it. Recorded so that it is not re-derived.
\end{remark}

\section{What the carrier cannot express}

\begin{proposition}[No reflexivity over element types]\label{prop:dag}
The requirement relation on $\El$, taking every alternative of every term, is a
directed acyclic graph without self-loops. Consequently no $X \subseteq \El$
contains a requirement cycle.
\end{proposition}

\begin{corollary}
Reflexive failure modes are not expressible over element types. They become
expressible over $\El \times \mathcal{A}$ for an asset set $\mathcal{A}$, under
the relation $\mathrm{backs} \coloneqq \mathrm{reads} \mathbin{;} \mathrm{over}^{-1}$.
\end{corollary}

\begin{conjecture}[Fibre separation]\label{conj:fibre}
There exist protocols $P \neq Q$ with identical element sets and materially
different solvency. Hence no function of $2^{\El}$ separates them, and separation
requires a party sort carrying obligor, attester and jurisdiction.
\end{conjecture}

\section{Composition}

\begin{definition}
$X \oplus Y \coloneqq \Gamma(X \cup Y)$.
\end{definition}

\begin{theorem}\label{thm:noncong}
$\oplus$ preserves observational equivalence but does not preserve
admissibility: there exist admissible $X, Y$ with $X \oplus Y \notin \Adm$.
\end{theorem}

\begin{remark}
The two claims are frequently conflated. Observational equivalence is defined by
quantification over all contexts, so its preservation is immediate; admissibility
is a different predicate and fails.
\end{remark}

\begin{measurement}\label{meas:frag}
A union-closed fragment $F \subseteq \Adm$ with $|F|/|\Adm| = 0.478$ certifies
$59\%$ of live protocol pairs with no errors. $F$ is a join-semilattice, not a
sublattice; maximality is open in $[11026,\,23055)$.
\end{measurement}

\begin{measurement}[Width dependence]\label{meas:width}
Composition success depends on maximum disjunctive width $w$: rates
$0.925, 0.710, 0.655, 0.684$ for $w = 1,\dots,4$, with union-closure falling from
$0.75$ to $0.13$ between $w=1$ and $w=2$. The observed $59\%$ lies in the
$w \in [2,4]$ band.
\end{measurement}

\begin{remark}[Corrected]
It is tempting to read Measurement~\ref{meas:width} as showing that the recurring
``good fragment'' phenomenon is a \emph{single} effect parameterised by
disjunctive width. That is half right, and the half that is wrong matters.
Proposition~\ref{prop:twoeffects} separates the two causes: prohibitions destroy
$\cup$-closure and nothing else, while disjunctive requirements destroy
$\cap$-closure and nothing else. Width governs one of the two failures, not both.
\end{remark}

\begin{proposition}[Two independent failures]\label{prop:twoeffects}
Let $\mathcal{M}(\Law)$ and $\mathcal{M}(\Haz)$ denote the model classes of the
requirements and the prohibitions. Then $\mathcal{M}(\Law)$ fails
$\cap$-closure exactly when some requirement has a term of width $\ge 2$, and
$\mathcal{M}(\Haz)$ fails $\cup$-closure exactly when $\Haz \neq \emptyset$. The
two failures are independent: neither implies nor mitigates the other.
\end{proposition}

\section{The structure theorem}

Commutation of $\Gamma$ and $\Delta$ would give a complete lattice of common
fixed points by Tarski's theorem, but commutation fails here. It can be weakened.

\begin{definition}[Invariance]
$\Delta$ is \emph{$\Gamma$-invariant} if $\Delta(\Fix(\Gamma)) \subseteq \Fix(\Gamma)$.
\end{definition}

\begin{theorem}[Invariance holds by construction]\label{thm:invariance}
If the warrant relation $C$ is the residual of the requirement relation $R$, then
$\Delta$ is $\Gamma$-invariant.
\end{theorem}

\begin{proof}
Let $X \in \Fix(\Gamma)$ and suppose $\Delta$ removes $e \in X$; then $e$ is a
dependent with $C(e) \cap X = \emptyset$. Suppose removal broke closure. Then
some $s \in X$ fires a requirement one of whose terms had $e$ as its only witness
in $X$. Since $C$ is the residual of $R$, any $s$ demanding a term containing $e$
is by construction a consumer of $e$, so $s \in C(e) \cap X$, contradicting
$C(e) \cap X = \emptyset$. Hence $\Delta(X) \in \Fix(\Gamma)$.
\end{proof}

\begin{corollary}[Structure]\label{cor:tarski}
$\Fix(\Gamma) \cap \Fix(\Delta)$ is a nonempty complete lattice.
\end{corollary}

\begin{proof}
$\Fix(\Gamma)$ is a complete lattice by Tarski's theorem. By
Theorem~\ref{thm:invariance} $\Delta$ restricts to it, and a monotone
self-map of a complete lattice has a complete lattice of fixed points, again by
Tarski.
\end{proof}

\begin{measurement}\label{meas:invariance}
Invariance was verified on the actual operators at $100\%$ in four independent
tests: the Galois closure $\gamma = \beta \circ \alpha$ over $58$ elements
($160/160$); the height-one definite closure $Cn$ over $58$ elements
($231{,}773$ distinct closures); $Cn$ \emph{exhaustively} over $2^{20}$
($160{,}000$ closed sets); and the scored admissibility predicate
($92{,}880/92{,}880$).
\end{measurement}

\begin{remark}[Scope, and a maintenance hazard]
Corollary~\ref{cor:tarski} governs $\Gamma \wedge \Delta$ only. Prohibitions are
downward-closed and lie outside it, so the non-congruence of composition
(Theorem~\ref{thm:noncong}) and the join-semilattice character of the composable
fragment (Measurement~\ref{meas:frag}) both stand. Tarski also gives a lattice,
not that meet is intersection.

More practically: Theorem~\ref{thm:invariance} holds \emph{because} $C$ was
derived as the residual of $R$. It is therefore a maintenance invariant, and it
breaks the moment a warrant row is hand-edited without its governing requirement.
\end{remark}

\section{Repair}

The prohibitions carry more structure than the requirements, and it is the
structure that answers the question a user actually asks: \emph{what is the
minimal change that restores admissibility?}

\begin{definition}[Clutter and blocker]
A \emph{clutter} on ground set $\El$ is an antichain $\mathcal{C} \subseteq 2^{\El}$.
Its \emph{blocker} $b(\mathcal{C})$ is the clutter of minimal sets meeting every
member of $\mathcal{C}$.
\end{definition}

\begin{remark}[Terminology]
In this tradition a \emph{transversal} meets each member exactly once; the object
we need is a \emph{cover}, meeting each member at least once. The distinction is
not cosmetic and the literature is not uniform about it.
\end{remark}

\begin{theorem}[Blocker duality; Isbell 1958, Edmonds--Fulkerson 1970]\label{thm:blocker}
$b(b(\mathcal{C})) = \mathcal{C}$.
\end{theorem}

\begin{corollary}[Minimal repair]
Let $X$ arm some prohibition. The minimal sets of elements whose removal restores
hazard-freedom are exactly the members of $b(\Haz)$ restricted to $X$.
\end{corollary}

\begin{proposition}[Addition and removal are one object]\label{prop:oneobject}
Minimal removals from a seed $S$ are given by the deletion minor
$b\bigl(\Haz \setminus (\El \setminus S)\bigr)$ and minimal forbidden additions by
the contraction minor $\Haz / S$. The two are exchanged by Seymour's identity
\[
  b(\mathcal{C} \setminus I / J) \;=\; b(\mathcal{C}) / I \setminus J .
\]
Hence a single computation answers both.
\end{proposition}

\begin{remark}
Theorem~\ref{thm:blocker} is evaluated purely combinatorially on the ground set:
it never evaluates $\emptyset$ and never iterates from $\bot$. It is therefore
\emph{immune} to Proposition~\ref{prop:collapse}, unlike every fixpoint framework
considered above. This is the first candidate structure that dodges the collapse
rather than being defeated by it.
\end{remark}

\begin{measurement}[The prohibition table is smaller than it appears]\label{meas:clutter}
Of the $20$ recorded prohibition rows, only $3$ are enforced positive element
sets, namely $\{Fl,Xm\}$, $\{Fl,Au,Rl\}$ and $\{Fl,Cp,Cl,Pl,Cd\}$. Nine name no
element at all; several are requirements written in negative form; and one pair
violates the antichain condition by containment. On the genuine $3$-member
clutter, $|b(\Haz)| = 9$, $b(b(\Haz)) = \Haz$ holds, the minor identities of
Proposition~\ref{prop:oneobject} hold on all eight elements involved, and the
maximal hazard-free sets are exactly the complements of $b(\Haz)$.
\end{measurement}

\begin{measurement}[Accessibility]\label{meas:access}
Every admissible set admits an admissible construction order from $\emptyset$:
verified $8{,}240/8{,}240$ across four $16$-element ground sets. Formally, for
every admissible $X \neq \emptyset$ there is $e \in X$ with $X \setminus \{e\}$
admissible.
\end{measurement}

\begin{remark}
Measurement~\ref{meas:access} is the ``seating order'' the vocabulary has always
implied and never recorded. Note that it does \emph{not} make $\Adm$ an
antimatroid: antimatroids require $\cup$-closure, which
Proposition~\ref{prop:twoeffects} says the prohibitions destroy. We obtain
accessibility without the axiom system that usually delivers it.
\end{remark}

\section{Open problems}

\begin{enumerate}
  \item \textbf{Invariance --- resolved} (Theorem~\ref{thm:invariance},
    Corollary~\ref{cor:tarski}). What remains is whether meet is intersection on
    $\Fix(\Gamma) \cap \Fix(\Delta)$, which Tarski does not supply.
  \item \textbf{Maximal composable fragment.} Characterise the largest
    $F \subseteq \Adm$ closed under $\oplus$; see Measurement~\ref{meas:frag}.
  \item \textbf{Minimal repair --- resolved for the prohibition half}
    (Theorem~\ref{thm:blocker}, Proposition~\ref{prop:oneobject}). What remains
    open is \emph{restoring closure after removal}, which is the completion
    problem and is NP-complete.
  \item \textbf{Independence.} Is the generating set $\El$ minimal, and which
    elements are definable from the rest?
\end{enumerate}

\end{document}
